In the wake of the various technological advancements made to securing firearm technology, it is imperative to understand the advantages that certain technologies bring and the potential risks they may pose. There have been a variety of approaches to firearm security, ranging from traditional lock-and-key mechanisms to biometric systems. The lock-and-key approach, while simple and effective, can be slow to unlock, which is not ideal in emergency situations. Biometric systems, on the other hand, offer quick access but can be unreliable in certain conditions, such as when the user has wet or dirty hands, and cannot be used by multiple users without reprogramming. Not to mention that biometric systems can be expensive and don't work with gloves. The Smart Trigger system, however, is designed to be a cost-effective and reliable solution that addresses the shortcomings of these traditional methods. It combines the speed of biometric systems with the reliability of traditional locks, while also being user-friendly and adaptable to the standard AR-15 platform.

Although, securing firearms is nothing new, our approach to firearm security is unique in its design and implementation. Instead of a costly biometric system or a traditional lock, the Smart Trigger system utilizes a low-cost, passive 125kHz RFID tag to provide quick and reliable access to the firearm. Embedded within a ring worn on the user's finger, this tag allows for immediate, reliable access without the need to fumble with keys or remove gloves. Using a unique RFID tag, the Smart Trigger system stays locked until the authorized tag is presented, allowing for quick access in emergency situations. Using a tag instead of a biometric system allows for multiple users to have access to the firearm without the need for reprogramming simply by distributing authorized rings, and it can be used with gloves or in adverse conditions. All processing electronics, including the RFID antenna and microcontroller, are housed within a custom 3D-printed grip, making the system compact, easy to install, and unobtrusive.

The system works as follows:
\begin{itemize}
    \item The user wears a ring with an embedded RFID tag on their finger.
    \item When the user grasps the firearm, the ring is naturally brought into range of the antenna embedded in the custom Smart Trigger grip.
    \item The internal system reads the RFID tag and verifies it against a list of authorized IDs.
    \item If the tag is authorized, a micro stepper motor inside the grip moves a blocking member away from the AR-15's native safety selector, allowing the user to switch the safety off and fire the weapon.
\end{itemize}

Once the safety is disengaged, it will not automatically re-engage, allowing the user to fire the weapon without any additional steps. To re-engage the safety, the user simply needs to switch the safety back on, which will re-engage the pin and lock the firearm again. 

This paper outlines the complete technical realization of the Smart Trigger system. It covers the System Architecture, Hardware Components and Integration, CAD Modeling and Design, Software Processes, and Related Work. The goal of this paper is to provide a comprehensive understanding of the Smart Trigger's capabilities as a supplementary, technical authorization step for the AR-15 rifle.